% ── §6b Methods: Evaluation Metrics ──
\section{Evaluation Metrics}
\label{sec:methods-metrics}

Proper evaluation of contact prediction methods requires careful attention to
the metrics used, the candidate universe over which they are computed, and the
baseline against which improvement is measured.  We describe each of these
elements here.

\subsection{Precision at Depth \texorpdfstring{$K$}{K}}
For a ranked list of candidate pairs and a set of native contacts, precision
at depth $K$ is defined as
\begin{equation}
	\mathrm{prec}@K = \frac{|\{\text{top-}K \text{ ranked pairs}\} \cap
		\{\text{native contacts}\}|}{K}.
	\label{eq:precision}
\end{equation}
We report precision at depths $K = L/10$, $L/5$, $L/2$, and $L$, where $L$ is
the protein length.  These depths follow the convention established by the
CASP contact prediction assessments~\cite{ezkurdia2009}.

\subsection{Recall}
Recall at depth $K$ measures the fraction of all native contacts recovered
within the top-$K$ predictions:
\begin{equation}
	\mathrm{recall}@K =
	\frac{|\{\text{top-}K\} \cap \{\text{contacts}\}|}{|\{\text{native contacts}\}|}.
	\label{eq:recall}
\end{equation}
Recall is bounded above by $\min(1, K / N_c)$ where $N_c$ is the total number
of native contacts; when $N_c > K$, perfect ranking still cannot achieve unit
recall.

\subsection{Area Under the ROC Curve}
The area under the receiver operating characteristic curve (AUC) provides a
threshold-free measure of ranking quality.  We compute AUC via the
Mann--Whitney $U$ statistic with mid-rank tie handling:
\begin{equation}
	\mathrm{AUC} = \frac{U}{n_+ \cdot n_-}, \quad
	U = R_+ - \frac{n_+(n_+ + 1)}{2},
	\label{eq:auc}
\end{equation}
where $R_+$ is the sum of ranks assigned to positive (contact) pairs, and
$n_+, n_-$ are the counts of positive and negative pairs respectively.  This
implementation handles tied scores correctly through average-rank assignment.

\subsection{Matthews Correlation Coefficient}
For binary classification at a fixed threshold, the Matthews Correlation
Coefficient (MCC) provides a balanced measure that accounts for all four
outcomes:
\begin{equation}
	\mathrm{MCC} =
	\frac{t_p \cdot t_n - f_p \cdot f_n}
	{\sqrt{(t_p{+}f_p)(t_p{+}f_n)(t_n{+}f_n)(t_n{+}f_p)}},
	\label{eq:mcc}
\end{equation}
which equals zero when either denominator factor vanishes.  MCC ranges from
$-1$ (perfect anti-classification) through 0 (random) to $+1$ (perfect).

\subsection{Enrichment Factor}
The enrichment factor expresses precision relative to the random-expectation
precision (the base rate):
\begin{equation}
	\mathrm{EF}@K = \frac{\mathrm{prec}@K}{N_c / N_{\mathrm{cand}}},
	\label{eq:enrichment}
\end{equation}
where $N_{\mathrm{cand}}$ is the total number of candidate pairs.  An
enrichment of $1\times$ means the method performs no better than random
ranking; values above $1\times$ indicate genuine signal.

\subsection{Baseline Clarification}
A common source of confusion in contact-prediction literature concerns the
meaning of ``random chance.''  Because contacts constitute only $\sim$3\% of
all candidate residue pairs, a method that ranks candidates randomly achieves
a precision equal to this base rate, not 50\%.  The value 0.5 would apply only
to a balanced binary classification problem where half the instances are
positive.  All enrichment factors reported in this work are computed relative
to the empirical base rate of each test set.
